\documentclass[a4paper,11pt]{article}

\usepackage[utf8]{inputenc}
\usepackage[T1]{fontenc}
% Utilisons palatino pour plaire à certains afficheurs de PDF
\usepackage{pxfonts} 
\usepackage{fancyhdr}
\usepackage{moreverb}
\usepackage{tocbibind}
\usepackage{geometry}
\usepackage{graphicx}
%\usepackage{rcs}
\usepackage{hyperref}
\usepackage[french]{babel}
\usepackage{listings}
\usepackage{url}
\usepackage{ulem}
\lstset{
    language=bash, %% Bon défaut dans notre cas
    inputencoding=utf8,
    literate=%
    {é}{{\'e}}1%
    {è}{{\`e}}1%
}
\usepackage{supertabular}

\usepackage{datetime2}
\usepackage[local]{gitinfo2}

% format de la date pour Git
\DTMnewdatestyle{mydate}{%
  \renewcommand*{\DTMdisplaydate}[4]{##3 \DTMfrenchMonthname{\DTMtwodigits{##2}} ##1}%
  \renewcommand{\DTMDisplaydate}{\DTMdisplaydate}%
}
\DTMsetdatestyle{mydate}

% definition du format de l'entete et du pied de page
\pagestyle{fancyplain}
\lhead{}
\rhead{\leftmark}
\rfoot{
  \begin{scriptsize}
  \textsl{version du \DTMusedate{gitdate}}
  \end{scriptsize}
}

%%% warning encadré
%%% utilisation : \begin{warning} blablabla \end{warning}
\usepackage{pifont,mdframed}
\newenvironment{warning}
  %%{\par\begin{mdframed}[linewidth=2pt,linecolor=red]%
  {\par\begin{mdframed}[linewidth=2pt]%
    \begin{list}{}{\leftmargin=1cm
                   \labelwidth=\leftmargin}\item[\Large\ding{43}]}
  {\end{list}\end{mdframed}\par}

%%% pour l'usage des programmes qui ne tiennent pas sur une ligne
%%% utilisation :
%%%%%% \minibox[frame]{
%%%%%%  \texttt{gengraphvect\.py [-h] -a ACQID [-o FORMAT] [-n NIVEAU]} \\
%%%%%%  \rightline{ \texttt{[-f FICHIER] [-s SENS] [-v] [VECTEUR ...]} }
%%%%%% }
\usepackage{minibox}


\hypersetup{
       colorlinks=true,%  Liens en couleur 
       linkcolor=blue,%   Liens bleus
       urlcolor=blue,%   Liens bleus
       %%hyperindex=true,%  Ajoute des liens dans les index.
       pdftitle={Documentation AIDA Volume 1 : documents de référence},
       pdfauthor={CNRM/GMEI/TRAMM - SAFIRE},
       pdfkeywords={Références AIDA}
}

\newenvironment{tab4col}{\begin{tabular}{|p{4cm}|p{2cm}|p{3.5cm}|p{5.5cm}|}}{\end{tabular}}
\newenvironment{tab3col}{\begin{tabular}{|p{4cm}|p{4cm}|p{7cm}|}}{\end{tabular}}
\newenvironment{tab2col}{\begin{tabular}{|p{4cm}|p{11cm}|}}{\end{tabular}}
\newcommand{\saute}{\vspace{0.6cm}}

% Let's have some nice headings
\addtolength{\headheight}{2.3pt}
\title{Documentation AIDA \\Volume 1 : documents de référence}
\author{CNRM/GMEI/TRAMM \and SAFIRE}
\date{Version du \DTMusedate{gitdate}}

%   Refmark sert à appeler plusieurs fois la même
% note de bas de page. Idée lue dans cahier Gut
%nr 15, avril 1993, pp 51-52
%%%\def\refmark#1{\hbox{\textsuperscript{\ref{#1}}}}
\usepackage[bottom]{footmisc}

\begin{document}
\maketitle
\newpage

\tableofcontents
\newpage

\section*{Présentation}
	AIDA est un \emph{atelier de traitement}
qui regroupe un ensemble de programmes utiles à l'analyse et
l'interprétation de séries temporelles de mesures de l'atmosphère
(issues de stations sol, de bateaux, d'avions, ou de tout autre plate-forme).
 Ces programmes sont interdépendants, et partagent un environnement
informatique commun. Ils doivent donc respecter certaines 
règles : formats de fichiers, nommages et organisation des
répertoires, variables d'environnement, etc. Ces règles sont divisées en
plusieurs chapîtres :
\begin{itemize}
  \item \nameref{part:arbo}, page \pageref{part:arbo}
  \item \nameref{part:formvec} ("\textit{séries temporelles}"), page \pageref{part:formvec}
  \item \nameref{part:meta}, page \pageref{part:meta}
\end{itemize}

% Partie : description de l'arborescence
\part{Répertoires et variables d'environnement}
\label{part:arbo}
\section{Répertoires}
\label{sec:repert}
   Tous les fichiers et programmes qui constituent AIDA se regroupent de
manière arborescente sous un unique répertoire racine. Il n'y a pas 
de contrainte sur le choix de l'emplacement de ce répertoire, qui peut 
varier d'une machine à l'autre.

 Dans ce répertoire-racine, se trouvent les sous-répertoires suivants :
\begin{description}
  \item[src] contient les sources des programmes qui constituent AIDA. Il 
   est divisés en sous-répertoires, un par langage (script, 
   python3, etc), plus un pour les bibliothèques générales écrites en C
   (\texttt{lib}) et un contenant les codes sources de la documentation
   AIDA (\texttt{doc}).

  \item[bin] contient tous les exécutables d'AIDA. Pour l'utilisation
   confortable d'AIDA, il est recommandé de rajouter ce répertoire
   dans sa variable PATH (voir \path{~/.profile.aida}) :

  \item[appli] contient le matériel de traitements des campagnes et des
    acquistions de l'utilisateur (scripts, programmes, fichires de
    configuration...). \\
    \textbf{Ce répertoire n'est pas géré par le gestionnaire de source de AIDA.}

  \item[meta] contient tous les fichiers de méta-données. Il contient au moins
   deux sous-répertoires utilisés par AIDA :
      \begin{description}
       \item[INF] : les meta-données des acquisitions sous forme de fichiers
                    <acqid>.inf.
       \item[campagnes] : les meta-données des campagnes sous forme de fichiers
		                <campagne>.camp.
      \end{description}
   L'utilisateur peut y ajouter les sous-répertoires utilisés par son matériel
   de traitement (par exemple, \texttt{crvol} pour les comptes-rendus de vol).
   Si le réseau local le permet, ce peut être un montage NFS ou un
   partage SMB.

  \item[local] contient les fichiers spécifiques à la configuration 
   informatique de la machine sur laquelle s'execute AIDA.

  \item[config] contient un certain nombre de fichiers 
   de configuration d'AIDA (liste des familles de paramètres, les définitions
   de graphiques).
   Ce répertoire peut être un montage NFS ou un partage SMB.

  \item[profils] contient un sous-répertoire spécifique pour chaque profil
   d'utilisation de AIDA (ex: 4M, EXEMPLES, SAFIRE, TRAMM).
   Chaque sous-répertoires contient au moins un fichier \texttt{plate-formes}
   (liste de plate-forme connues pour ce profil d'utilisation de AIDA).

  \item[share] contient des fichiers partageables (cartes, logos...).

  \item[data] contient toutes les données. Pour des raisons pratiques,
   il s'agit souvent d'un lien vers un gros disque dédié aux données.
   Si le réseau local le permet, ce peut être un montage NFS ou un
   partage SMB.
   L'organisation des données est décrite un peu plus loin 
   (paragraphe \ref{sec:arbodata}).

  \item[doc] contient la documentation d'AIDA (donc, en particulier, ce
   fichier). Attention, ce n'est pas la même chose que la documentation
   sur les vols ou les traitements, produite par AIDA.

  \item[include] : contient tous les fichiers include (\path{*.h} du C, ...).

  \item[lib] contient toutes les bibliothèques AIDA (fichiers \path{*.a}, 
   \path{*.sl}, \path{*.ln}, \path{*.so} mais aussi sous le répertoire
   \path{pythonX.Y}, les librairies python de AIDA).

  \item[log] pour les programmes qui veulent laisser une trace de ce qui
   s'est passé pendant leur éxécution.

  \item[extern] accueille le matériel tierces utilisé par AIDA, lorsqu'elles
   ne sont pas fournies par le système (par exemple \texttt{cextract}, ...)
\end{description}
 
\section{Variables d'environnement}
   L'emplacement du répertoire racine pouvant varier, on utilise une 
variable d'environnement, \textsf{AIDA\_HOME}, pour en 
désigner l'emplacement. Tous les programmes prennent en compte cette
variable, il est donc indispensable qu'elle soit établie dans 
l'environnement de chaque utilisateur.

   A l'installation de AIDA, les variables d'environnement suivantes sont
déclarées (voir \path{~/.profile.aida}) :
\begin{description}
  \item[\textsf{AIDA\_HOME}] désigne l'emplacement du répertoire racine de AIDA.
  \item[\textsf{AIDA\_PROFIL}] désigne le profil d'utilisateur AIDA.
  \item[\textsf{AIDA\_META}] désigne l'emplacement du répertoire des
   méta-données.
  \item[\textsf{AIDA\_DATA}] désigne l'emplacement du répertoire des données
\end{description}

\medbreak
   D'autre variables peuvent désigner chacun des sous-répertoires
décrits au paragraphe \ref{sec:repert}. Elle sont utiles lorsqu'un
utilisateur ne souhaite pas utiliser l'arborescence commune mise en place,
par exemple pour accéder à des données situées ailleurs que dans
\texttt{\$AIDA\_HOME\//data}, ou à des versions spécifiques
de programmes situés ailleurs que dans \texttt{\$AIDA\_HOME\//bin}.\\
La règle pour construire le nom de ces variables est simple : elles sont
baptisées \textsf{AIDA\_\textit{xxx}}, où \textit{xxx} est le nom en majuscule
du sous-répertoire qu'elle désigne. \\
Ainsi :
\begin{itemize}
  \item[\texttt{AIDA\_BIN}] : sert à accéder à des programmes situé ailleurs
    que dans \texttt{\$AIDA\_HOME\//bin},
  \item[\texttt{AIDA\_LIB}] : sert à utiliser des librairies situées ailleurs
    que dans \texttt{\$AIDA\_HOME\//lib},
  \item[\texttt{AIDA\_DATA}] : sert à travailler sur des données situées
   ailleurs que dans \texttt{\$AIDA\_HOME\//data}.
\end{itemize}

\section{Identification de chaque période de mesure}
\label{sec:AcqId}
   On appelle \textit{période de mesure} le temps pendant lequel
un système d'acquisition a fonctionné. Pour les systèmes qui
fonctionnent par période, la définition est assez simple (un 
lâcher de ballon-sonde, un vol d'un avion). Pour les systèmes qui
fonctionnent en continu, la période de mesure est en général
définie comme une journée, mais rien n'empèche  de choisir une
durée plus courte ou plus longue (le mois, pas exemple).

\medbreak
On distingue 5 types de période de mesure :

\textbf{- Journee} (\textit{anciennement Acquisition})

\textbf{- Vol}

\textbf{- Lacher}

\textbf{- Lacher\_en\_vol}

\textbf{- Drone}


\medbreak
   Pour des raisons pratiques (nommage de fichiers, choix d'une
période parmis les autres, ..), il est important de pouvoir désigner
chaque période de mesure par un nom unique, appelé \textbf{\texttt{acqid}}.

\medbreak
Un \textbf{\texttt{acqid}} est une chaîne de 8 caractères, de la forme : \\
\begin{center}
{\fbox{ \texttt{ppAAnnnN}}}
\end{center}

\medbreak
\begin{itemize}
  \item Les 2 premiers (\texttt{pp}) sont alpha-numériques, et désignent la 
   plate-forme de mesure. La correspondance entre cette abbréviation et
   le nom usuel, la description, etc se fait par le fichier
   {\texttt plate-formes} décrit au paragraphe \ref{sec:platef}, page 
   \pageref{sec:platef}.

  \medbreak
  \item Les 2 suivants (\texttt{AA}) sont numériques, et désignent l'année.
   Les valeurs supérieures à 70 correspondent aux années 1970 à 1999,
   et les autres aux années 2000 à 2069.

  \medbreak
  \item Les 4 suivants (\texttt{nnnN}) désignent le
   numéro d'acquisition de manière différente selon le type de période de
   mesure :

   \begin{itemize}

     \item[\textbf{Journee}] : \texttt{nnnN} est le quantième du jour dans
                               l'année (numérique)

     \item[\textbf{Vol}] : \texttt{nnnN} est le quantième du vol dans l'année
                           (numérique, sur 4 chiffres commençant par 0)

     \item[\textbf{Lacher}] : \texttt{nnnN} est le nombre d'heures depuis le 01/01
                              de l'année (numérique)

     \item[\textbf{Lacher\_en\_vol}] : \texttt{nnn} est le numéro du vol (numérique),
                                       \texttt{N} est le numéro du lâcher
                                       (alphanumérique [0-9] puis [a-z] puis [A-Z]
                                       soit 62 numéros de lâcher)

     \item[\textbf{Drone}] : \texttt{nnn} est le quantième du jour dans l'année
                             (numérique), \texttt{N} est le numéro du vol (idem
                             numéro de lâcher pour \texttt{Lacher\_en\_vol})
   \end{itemize}
\end{itemize}

\textbf{\textit{Note : avant 2008, la partie \texttt{nnnN} des acqids de type
vol n'était constitué que de 2 chiffres.}}


\section{Détail des arborescences principales}

Certains des sous-répertoires présentés ci-dessus contiennent
une information volumineuse et/ou complexe : ils sont en fait le sommet
d'une arborescence. Nous décrivons ici les plus importantes.

\subsection{L'arborescence des données} 
\label{sec:arbodata}

    Dans le cadre d'un fonctionnement normal, l'utilisateur 
n'a pas besoin de connaître les détails du stockage
des données.

    Les programmes (et les bibliothèques sous-jacentes) sont conçues 
pour qu'il suffise à l'utilisateur de spécifier 
l'identification d'une période de mesure (par exempl
 l'\textbf{\texttt{acqid}} \texttt{Tf210356})
pour que les programmes d'AIDA sachent accéder aux données.

   Pour que cela fonctionne, il faut évidement spécifier une 
organisation de stockage, qui doit être respectée par tous les 
programmes.  Actuellement, cette organisation du stockage repose
sur le système de fichiers et de répertoires arborescents 
fourni par le système d'exploitation. C'est une arborescence 
à trois niveaux : le \textbf{premier correspond à la plate-forme,
le second à l'année, et le troisième à la période de mesure}.

    Cette organisation permet d'avoir une architecture complètement
ouverte, en permettant au plus grand nombre de logiciels
et de langages des programmation de travailler sur les données d'AIDA.

    L'emplacement d'un fichier se déduit automatiquement
du nom unique donné à chaque période de mesure :

    Le répertoire \texttt{\$AIDA\_HOME\//data} est le
point de départ de toute l'arborescence des données. Il contient un
répertoire par plate-forme de mesure. Le nom de ces répertoires
est exactement le même que celui inscrit dans le fichier
\texttt{plate-forme} (décrit au paragraphe \ref{sec:platef}).
A son tour, chacun de ces répertoires contient un répertoire par 
année, dont le nom est écrit en quatre chiffres, et qui contient
un répertoire par période de mesure. Enfin, dans ce répertoire
«feuille», il y a un fichier par
paramètre, qui contient la série complète des données (format
\textit{vecteur} décrit en détail dans la partie suivante).

 En reprennant l'exemple de  l'acquisition \texttt{Tf210356}, le 
répertoire sera : \\
\texttt{\$AIDA\_HOME\//data\//toulouse-flux\//2021\//0356}.

\subsection{L'arborescence des applications} 
    Pour chaque campagne de mesure, il est nécessaire
d'avoir des fichiers particularisés : des \textit{scripts}
de traitements, des scripts de production des fichiers
ASCII, ....

     Ces fichiers correspondent à des \textbf{applications}
utilisant AIDA, et sont regroupés dans \texttt{\$AIDA\_HOME\//appli},
plus précisément dans :
\texttt{\$AIDA\_HOME\//appli\//script\//\$nom\_campagne}.




% Partie : Format "vecteur"

% Partie 2 : Format des fichiers vecteurs
\part{Format des fichiers vecteurs}
\label{part:formvec}
\section{Présentation générale}
  Les \emph{fichiers-vecteurs} sont destinés au stockage de séries
temporelles de paramètres mesurés ou calculés. Ils associent
étroitement la donnée (dans un format binaire,
copie directe de la représentation en mémoire de
l'ordinateur), et sa description (dans une en-tête au format ASCII).
%Ce format d

  L'en-tête descriptive contient toutes les informations
utiles (format, unité, appareil ou algorithme d'origine...), \textbf{sauf}
la datation, car celle-ci
est parfois très complexe à gérer (variables
asynchrones). A chaque fichier de données, on associe donc un fichier
temps qui indique l'instant de la mesure de chaque
échantillon. Il y a strictement concordance entre ces deux fichiers :
la date de la i\up{ème} donnée est la i\up{ème}
valeur du fichier temps.

      Naturellement, plusieurs fichiers de données
synchrones sont associés au même fichier \og temps\fg{} : on
ne duplique pas inutilement ces fichiers (dans le cas simplifié
où toutes les variables seraient acquises à la même
fréquence par le même système, il n'y aurait qu'un seul vecteur-temps).

\section{Contenu}

Ces fichiers sont organisés en 2 parties :
\begin{enumerate}
  \item Un en-tête descriptif (souvent nommé \og header\fg{})
  \item Les données en binaire.
\end{enumerate}

\subsection{Format de l'en-tête.}
L'en-tête est un bloc d'une taille multiple de 512
octets, qui contient quelques informations techniques et le descriptif de la
variable dans un bloc ASCII. Ce bloc ASCII est encadré entre deux
marqueurs de début et de fin.
Voici la structure précise de l'en-tête :
\begin{itemize}
  \item Une chaîne de repérage : \og VECT-DATA\texttt{\textbackslash{}n}\fg{}
  \item La taille de l'en-tête (complète, avec les marqueurs, le bloc
ASCII, le remplissage de fin), en octets. Cette taille est forcément un
multiple de 512, et elle est écrite en ASCII cadré à
droite sur 10 caractères suivis d'un saut de ligne 
(\og \texttt{\textbackslash{}n}\fg{}).
  \item le marqueur de début du bloc ASCII, qui se compose 
    d'un caractère \og STX\fg{} (\texttt{\textbackslash{}002} ou \^{}B)
    et d'un caractère \og LF\fg{} (\texttt{\textbackslash{}n}).
  \item Le bloc ASCII, portion de texte documentant la variable, dont la
description précise est donnée plus loin.
  \item Le marqueur de fin du bloc ASCII, qui se compose de :
  \begin{itemize}
    \item une chaîne \og FIN\_HDR\texttt{\textbackslash{}{}n}\fg{}
    \item la suite des caractères \og NP\fg{} (\texttt{\textbackslash{}}0cx ou \^{}L),
       \og ETX\fg{} (\texttt{\textbackslash{}}003 ou \^{}C),
      \og SUB\fg{} (\texttt{\textbackslash{}}1ax ou \^{}Z)
  \end{itemize}
\end{itemize}
Ensuite, cet en-tête, jusqu'à la fin (c.a.d jusqu'au prochain
multiple de 512 octets), est rempli de caractères \og \fg{} (espace).


\subsubsection{Syntaxe du \og bloc-ASCII\fg{}}
Le texte décrivant la variable prend la forme suivante :
\begin{itemize}
  \item une première ligne contenant le nom de la variable (normalement
identique au \texttt{basename} du nom de fichier), encadré de crochets.
  \item les lignes suivantes sous la forme \og \textsf{attribut = valeur}\fg{}
\end{itemize}
Les caractères sont codés selon la norme ASCII, ou
éventuellement selon son extension ISO-8859-1 (ou ISO-Latin1).
Le caractère de fin de ligne est le \og NL\fg{}
(Il n'y a pas actuellement de spécification de longueur de ligne maximale).

%%%\paragraph{attributs obligatoires}
Les attributs suivants sont obligatoires :

\begin{description}
  \item[Nom du paragraphe] (entre crochet) nom du paramètre
et \texttt{basename} du nom de fichier.
  \item[vol] nom unique attribué à la série de mesure. C'est une chaîne de
6 à 8 caractère au format \texttt{PPaann[nn]}, et qui est décrite de façon
détaillée section \ref{sec:AcqId} page \pageref{sec:AcqId}. Pour des raisons
de compatibilité historique, ce champ s'appelle exactement «vol», même pour des
mesures de station de surface, ou à bord de bateaux.
  \item[format] le codage informatique des données. La liste des
  formats connus par la version actuelle d'AIDA est la suivante :

  \begin{description}
    \item[\textsf{Byte}, \textsf{Int\_16}, \textsf{Int\_32}, \textsf{Int\_64}] :\\
       entiers signés
    \item[\textsf{U\_Int\_8}, \textsf{U\_Int\_16}, \textsf{U\_Int\_32}] : \\
       entiers non signés
    \item[\textsf{Ieeefloat\_32}, \textsf{Ieeefloat\_64}] : \\
       Nombres \og flottants\fg{} codés selon la norme ANSI/IEEE 754-1985.
    \item[\textsf{Quanta\_12\_bits},
          \textsf{Quanta\_14\_bits},
          \textsf{Quanta\_16\_bits},
          \textsf{Quanta\_32\_bits}] : \\
       «mots-machine» provenant d'un convertisseur analogique-numérique. \\
       Les \textsf{Quanta\_12\_bits} et \textsf{Quanta\_14\_bits} sont
       stockés dans des mots de 16 bits, avec seulement 12 (ou 14) bits
       significatifs.
    \item[\textsf{Quanta\_14\_bits\_signe}] : \\
       la version signée d'un \textsf{Quanta\_14\_bits}.
    \item[\textsf{Arinc\_SY429}] : \\
       codage selon la norme aéronautique ARINC-SY429 (mot de 64 bits)
    \item[\textsf{BC635\_Time}] : Codage produits par la carte BC635.
    \item[\textsf{Special}] : Codages particuliers.
    \item[\textsf{Hpfloat\_32}] : \\
       nombre «flottantsi» codés à la manières des machines HP-1000.
    \item[\textsf{NCAR2D\_time}] : Le champ de datation des échantillons 
«2D» enregistrés par la carte développée au NCAR.
  \end{description}

  Dans les noms de formats, la casse (majuscule/minuscule) n'est pas signifiante.
  \item[endianness] indique l'ordre de la séquences d'octets des données. Ce champ peut prendre la valeur \texttt{little} ou \texttt{big}. Ce champ a été ajouté en 2021. Afin de conserver la compatibilité avec les anciens vecteurs, si ce champ n'est pas présent alors on considère endianness = big. 
  \item[dimension] taille en octet d'une
représentation de la variable. Si la valeur de l'attribut
\textsf{format} était un type connu du
langage C, on aurait : \textsf {dimension = sizeof(format)}. En
général, donc, \textsf{dimension} peut se
déduire de \textsf{format}, l'exception étant bien sûr le format 
\textsf{Special}.
  \item[dim\_array]  Dans le cas où chaque échantillon de la
variable est représenté par un tableau au lieu d'une simple
valeur (par exemple les spectres des FSSP), on indique ici les dimensions de
ce tableau. Dans le cas où le tableau a plus d'une dimension (une
image, par exemple), les tailles sont indiquées séparées par un \textsf{|}.
L'ordre des tailles indiquées suit la convention
du langage C. Ce champ n'est pas obligatoire; en son absence on suppose
que la variables est scalaire (un seul échantillon par instant de mesure).
   \item[ref\_time] Nom (\og logique\fg{}, sans le \texttt{.vbn} à la fin)
du paramètre qui contient la datation de cette
variable; ce champ n'est généralement pas présent
dans les variables qui contiennent elle-même un temps, mais il n'est pas
vraiment interdit qu'il y soit présent, dans les cas où l'on
souhaiterais comparer plusieurs systèmes, et donc savoir à
quelle heure (du premier système) correspond l'heure (du deuxième).
   \item[time\_bnds] dans ref\_time, précise les bornes temps de 2 manières possibles : 
   
time\_bnds = vector : <nom du vecteur de bornes temps> 

time\_bnds = decal\_av, decal\_ap (exemple : time\_bnds = -1800, 0 ) 

les bornes temps sont calculées par time+decal\_av, time+decal\_ap 

le nom du vecteur bornes temps est sous la forme : rft\_bnds\_<chaine libre>\_freq
    \item[datation\_convention] attribut informatif dans le reftime pour indiquer la norme de codage
     
	datation\_convention = N avec N variant entre 0 et 1  (0, : debut, 1 : fin intervalle)
    \item[type\_ech] Indique la façon dont la variable a
   été échantillonnée. Deux valeurs sont 
  possibles : \textsf{regulier} et \textsf{irregulier}.
   Ce champ n'est obligatoire que pour les vecteurs de datation; 
  S'il n'est pas présent dans un vecteur de données, on le déduit du
vecteur de datation (celui désigné par
\textsf{ref\_time}); et si celui-ci est absent, on finit par supposer que
l`échantillonnage est irrégulier.
    \item[freq\_ech] Si la variable a été
régulièrement échantillonnée, la valeur de la
fréquence d'échantillonnage, en Hz\footnote{Attention, cette
fréquence n'est pas forcément un entier $\ge$ 1, c'est un
réel $>$ 0}. Cette valeur peut être écrite avec toutes les décimales
nécessaires (dans une limite raisonnable, disons 30 caractères), mais
les expressions mathématiques (par exemple $ 1 / 3$) ne sont pas admises.
Pour les variables irrégulièrement échantillonnées (\textsf{type\_ech = irregulier}), ce champ est
présent et vaut 0.\\
  (Pour préserver une compatibilité
avec les précédentes versions du format, qui utilisaient le
mot-clef \textsf{sample}, les programmes qui lisent des fichiers vecteurs
doivent prévoir le cas où \textsf{freq\_ech} est absent, et
lire alors le champ \textsf{sample}).
    \item[unit] une unité répertoriée par la librairie \textit{udunits} du
                groupe unidata ou \textsf{unitless} (pour les valeurs issues
                d'un système de numérisation, l'unité la plus appropriée est
                \textsf{bit}). \\
                Pour des raisons historiques, on peut trouver l'unité
                \textsf{without} qu'il convient alors d'interpréter comme
                \textsf{unitless}.
%%%%On trouvera en annexe 1 une liste
%%%%des unités les plus fréquemment employées.
  \item[orientation] lorsque le paramètre contenu dans un vecteur est
une grandeur physique orientée (un angle, une distance sur un axe, une
composante d'un vecteur), ce champ contient une description de cette
orientation. Par exemple, pour un roulis : «~>0 lorsque le coté droit
de l'avion (ou du navire) est enfoncé~». La présence de cet attribut n'empêche
pas d'utiliser un nom explicite, lorsque c'est possible, par exemple
"\texttt{ray\_glo\_d\_...}" pour un rayonnement global descendant, ou ``\texttt{vit\_e\_\ldots}"
pour une composante de vitesse positive vers l'est.

  \item[date\_creation] au format \textsf{jj/mm/aaaa hh\string:mm\string:ss}.
  \item[origine] est un champ qui décrit à partir de quoi la variable a été
    crée.  Il est composé de plusieurs mots séparés par des espaces. Le
    premier décrit la fonction du programme producteur : c'est soit
    \textsf{lecture}, soit \textsf{calibration}, soit \textsf{synchronisation}, 
    soit enfin \textsf{traitement}.\\

    La suite du champs précise la description. Elle varie selon les cas :

    \begin{description}
      \item{Dans le cas \textsf{lecture}}, le mot suivant est le type de
        fichier d'origine [bande/f\_turb/NetCDF,...] puis le nom (sans le
        chemin complet) si il s'agit d'un fichier.

      \item{Dans le cas \textsf{calibration}}, le mot suivant est le nom du
        fichier vecteur d'origine, suivi de l'éventuelle méthode employée pour
        la calibration.

      \item{Dans le cas \textsf{traitement}}, la suite indique le  nom du
        programme de traitement (par compatibilité historique, si la suite est
        vide, on considérera que le programme d'origine est adc).
    \end{description}

\end{description}

%%%\paragraph{attributs optionnels}
  Les champs suivants sont optionnels, mais il est recommandé de les
renseigner chaque fois que c'est possible :
\begin{description}
   \item[name] Nom complet du paramètre, ne dépassant pas une
ligne, en français.
   \item[english\_name] Nom complet du paramètre, ne dépassant 
pas une ligne, en anglais (dans un fichier NetCDF, ce serait l'attribut
\textsf{long\_name}).
   \item[information] Description textuelle élargie, pouvant compléter 
le champ «name». 
  \item[date\_vol] la date de mise en route de l'acquisition, au 
format \textsf{jj/mm/aa hh\string:mm\string:ss}. L'information
sur le jour doit correspondre au champ date de la section [IDENTIFICATION]
du fichier .inf de l'acquisition. Si ces informations sont contradictoires,
le fichier .inf fait foi.
%% Ajouts suite à la réunion CAM/TRAMM du 02/12/1999
   \item[bornes] Dans le cas où chaque échantillon du paramètre est un 
tableau (attribut \textsf{dim\_array} supérieur à 1), ce champ précise 
les bornes de chacun des éléments du tableau.
Les bornes peuvent être être définie littéralement, sous la forme d'une expression permettant de les
calculer, ou être contenues dans un autre vecteur pour le cas où les bornes sont variables dans le temps.
La syntaxe peut prendre 3 formes :
\begin{itemize}
\item une suite de couples sous la forme :\\
\texttt{bornes = ($x_1$,$y_1$); ...; ($x_n$,$y_n$) }.\\
Les $x_i$ sont en général des réels, ou le caractère - si la
limite n'est pas définie (+ ou - l'infini). Le nombre de
couples est la valeur de \textsf{dim\_array}. (\textbf{Exception} : Les
vecteurs "quanta" provenant de la carte-coupleur \textsc{CNRM-1D} contiennent
une structure complexe : tableau de comptage et voies annexes. Dans 
ce cas, le champ \textsf{bornes} ne décrit que le tableau de comptage,
et le nombre de couples correspond à la valeur de l'attribut \textsf{classes}).

\item une formule permettant de calculer les bornes : \\
bornes = <range> : <debut> $d$, <fin>$f$, <nclasses>$N$, <pas>$P$, <largeur>$L$
Les valeurs des champs $d$, $f$, $N$, $P$, $L$ sont des valeurs numériques réelles.\\
L'ordre des champs n'a pas d'importance et 3 d'entre eux sont obligatoires.\\
Le champ <pas> est obligatoire dans tous les cas.\\
Le champ <largeur> est optionnel.  S'il n’est pas présent, alors $L = P$ \\
Parmi les 3 champs : <debut>, <fin> , <nclasses>, 2 d’entre eux sont obligatoires \\
La relation entre les champs est :  $f = d + (N -1)*P + L$ 

\item la réference à un  autre vecteur contenant les bornes variables dans le temps : \\
<vector> : \textsf{vecteur\_bornes}. Ce vecteur a 2 dimensions \textsf{dim\_array} et 2. Il contient à chaque pas de temps les bornes de début et fin de chaque classe.\\ 

\end{itemize}

Pour les vecteurs à plusieurs dimensions, les bornes de chaque dimension sont décrites par 
les mêmes syntaxes séparées par des \textsf{|} si plusieurs dimensions. \\
Par exemple : 

\-\hspace{1cm} {bornes = <range>: .... | <vector>: ...  | ($x_1$,$y_1$); ...; ($x_n$,$y_n$)}

L'ordre suit la même règle que \textsf{dim\_array}.

   \item[unit\_bornes] Indique l'unité de l'attribut \textsf{bornes}
décrit ci-dessus, séparés par des \textsf{|} s'il y a plusieurs dimensions.
   \item[type\_acq]indique la façon dont la variable a
    été acquise. Deux valeurs sont possibles :\\
    \textsf{synchrone} (dans le cas où l'instant de l'acquisition est
     commandé par le système d'acquisition),\\
     ou \textsf{asynchrone}
      (dans le cas où l'appareil lui-même décide du moment où
      il envoie ses données au système d'acquisition).
   \item[position] Un mot qui désigne un emplacement sur la plate-forme
    de mesure.

   \item[misval] la valeur spéciale pour les cas où la valeur
           n'est pas disponible.
   \item[invval] la valeur spéciale pour les cas où la valeur
           était manifestement hors-norme.
   \item[obs\_min] La plus petite valeur contenue dans la série
   \item[obs\_max] La plus grande valeur contenue dans la série
   \item[normal\_min] La plus petite valeur physiquement possible.
   \item[normal\_max] La plus grande valeur physiquement possible.
   \item[famille] Le type de grandeur physique de la mesure. La liste des
familles connues est contenue dans le fichier \path{famille}, décrit
en section \ref{sec:familles}, page \pageref{sec:familles}.
   \item[fct\_origine] Une description de la fonction de calcul qui a
produit ce vecteur.  Dans les cas où l'origine est «~lecture~» ou
«~calibration~», il peut servir à donner des précisions sur le calcul
effectué (formule, valeur des coefficients). Dans le cas où l'origine
est «~traitement~», il doit indiquer
tous les éléments d'entrée du calcul (vecteurs, coefficients), ainsi
que les éléments de sortie (vecteurs ou paramètres produits
simultanément). Dans le cas de traitement par adc, ce champ est une copie
de la ligne de commande du calcul.
   \item[fct\_comment] Commentaires ou compte-rendu produit(s) par la
fonction de traitement.
   \item[origine\-WARNING] messages d'avertissement ou de détection
d'anomalie qui ont été émis lors de la fabrication du fichier-vecteur
(par exemple « unité non reconnue »).

\end{description}

  Pour les vecteurs issus directement d'instruments (les \og quanta\fg{}
et \og cal\fg{} ), il existe deux méthodes (deux séries de champs optionnels)
pour décrire la mesure. Les programmes recherchant cette information doivent
donc savoir traiter les deux cas.

Le premier modèle est simple, et se base sur trois attributs :
\begin{description}
   \item[sensor] : nom de l'appareil de mesure.
%% ajouts 2005
La désignation complète d'un instrument comporte : le nom du constructeur,
le nom du modèle, le numéro de série, et une éventuelle référence interne 
(inventaire, ...).
Pour un meilleur traitement automatique ultérieur, il est recommandé de  
fournir ces quatre informations, séparées par des ":".
%% Fin ajouts 2005
   \item[transducer] : transducteur. Format exact non spécifié.
   \item[calibration] : méthode et date de calibrage ou étalonnage. Format exact
non spécifié.
\end{description}

Le second modèle est plus complexe, mais plus détaillé :
\begin{description}
%% Fin ajouts suite à la réunion CAM/TRAMM du 02/12/1999
   \item[chaine\_acquisition] (anciennement chaine\_mesure obsolète) : pour les paramètres bruts, description de la chaîne 
d'acquisition, selon le format : \\
\texttt{chaine\_acquisition = module\_1 module\_2 module\_3 module\_4 ...}

Les noms des modules sont de la forme module\_N, toujours en minuscules,
et séparés par des espaces.

   \item[module\_N] : description du N\up{ieme} module de la chaîne
     d'acquisition, suivant le format : \\
     \texttt{module\_N = <type\_module>:<fabricant>:<Modèle>:<Numéro de série>:<Numéro de référence interne>:<date et lieu de calibrage ou étalonnage>}

\smallbreak
Il n'est pas interdit qu'il y ait plusieurs modules de même
type dans une chaîne d'acquisition.\\

\smallbreak
La définition précise des types est : \\
  \begin{tabular}{rp{10cm}}
    antenna & Élément mécanique permettant de capter
    soit le flux d'air et de l'amener sur l'élément sensible,
    soit un signal radio-électrique (radars, récepteurs GPS, ...) \\
    
    sensor & Élément sensible. Une de ses propriétés physique varie en fonction
    de la grandeur que l'on veut mesurer. \\
    
    transducer & Transforme une propriété physique en tension, fréquence, 
    courant, ... Il peut être constitué d'un ensemble sensor+transducer 
    lorsqu'ils sont fortement liés (sensor pas facilement interchangeable) \\
    
    conditionner & Circuit électronique de traitement du signal issu 
    d'un capteur ou d'un transducteur. Il est utilisé pour simplifier l'interface
    du capteur avec le système d'acquisition de données. Il intervient entre le 
    capteur ou le transducteur, et l'unité d'acquisition. \\
    
    acquisition\_unit & Interface entre le monde extérieur (tensions, bus 
    numériques, lignes RS232) et le système d'acquisition.\\
    
    instrument & Appareil autonome se suffisant à lui-même, disposant 
    d'un affichage ou d'un système de stockage des données.\\
  \end{tabular}
\end{description}

%%%\paragraph{attributs libres}
  Enfin, tous les noms qui ne sont pas décrits ici sont libres d'utilisation
pour certaines applications spécifiques. En particulier les vecteurs résultant
d'une simple lecture, ou d'une calibration, conservent, pour que
l'auto-documentation soit complète, les attributs qui sont définis dans le
fichier d'origine (en-tête de fichier NetCDF, fiches de configuration du
système d'acquisition du CAM, etc...).

  Le tableau suivant recense les mots déjà définis par le système
d'acquisition du CAM, que d'autres applications doivent donc éviter
d'utiliser : \\

  \begin{tabular}{p{5cm}p{5cm}p{5cm}}
      bloc                 & calibration                 & channel      \\
      classes\footnotemark[1]
                           & cnv\texttt{*} (cnvi, cnvj, cnvrtn, ...) & codage \\
      code                 & codeRT                      & coupleur     \\
      cutoff\_freq         & delay                       & gain         \\
      index                & label                       & law          \\
      laws                 & mode                        & module       \\
      msgindex             & range\footnotemark[2]
                                                         & raz          \\
      recup                & reftime\footnotemark[3]     & resolution   \\
      sample\footnotemark[4]
                           & transducer                  & typeofdata   \\
      version              & voiesup\footnotemark[1]     & voie\_sup\_\textit{x}\\
  \end{tabular}

%%% notes de bas de page
\footnotetext[1]{\textsf{classes} et \textsf{voiesup} servent à préciser
  \textsf{dim\_array} pour les vecteurs «quanta» de sondes 1D. Normalement,
  \textsf{dim\_array}=\textsf{classes} + \textsf{voiesup}.}
\footnotetext[2]{Attention : cet attribut n'a pas le même sens pour les données
  issues de l'INS et des cartes analogiques.}
\footnotetext[3]{Attention : ne pas confondre avec ref\_time !!!}
\footnotetext[4]{Attention : n'est pas obligatoirement exprimée en Hz ! Voir
  aussi le mot-clef «blocage» dans les descripteurs de blocs}  

\subsection{Format de la partie binaire}
Cette partie du fichier contient la suite des
échantillons de la variable, dans le format indiqué dans la
partie descriptive.

    Le cas des variables qui ne sont pas de simples scalaires, mais des
tableaux (\texttt{dim\_array>1}) nécessite quelques précisions.  Le principe
général est que chaque tableau est considéré comme un échantillon, et le fichier
contient donc la suite chronologique des tableaux. Lorsqu'un tel fichier sera
chargé directement en mémoire, la façon de le décrire sera donc différente
selon le langage informatique utilisé.

Par exemple, pour des données se présentant sous la forme d'un tableau 
unidimensionnel (dim\_array=n), la déclaration correspondante en C sera :\\
\texttt{float data [nombre\_de\_secondes] [taille\_du\_tableau]}, ou encore\\
\texttt{float data [] [taille\_du\_tableau],}

alors qu'en FORTRAN, on écrirait plutôt :\\
\texttt{DIMENSION data (taille\_du\_tableau, nombre\_de\_secondes)}, ou encore\\
\texttt{DIMENSION data (taille\_du\_tableau, *)}

Une autre façon de voir les choses est de considérer un programme "universel", qui 
doit être capable de traiter des variables dont le dimensionnement est différent;
il faut alors considérer la zone de données comme un tableau unidimensionnel, 
et faire lui-même les calculs d'indice. La formule qui donne l'indice $k$ de la $i$\up{ème}
valeur de la $s$\up{ème} seconde de donnée est : $k = i + s*\textrm{taille\_du\_tableau}$.

Les cas de nombre de dimensions supérieures (par exemple les images pour lesquelles
\textsf{dim\_array=m|n}) se traitent de la même façon;
la formule générale qui donne l'indice $k$ du pixel $i$,$j$ de
l'image de la $s$\up{ème} seconde est :
$ k = s*m*n + i*m + j $

\subsection{Cas particulier des variables temps}

Elles contiennent toutes un temps depuis 00h00 UTC le jour de la mise en route du
système d'acquisition, exprimé 
en secondes. (Attention, pour les acquisitions qui chevauchent minuit
UTC, le nombre de seconde peut donc dépasser 86400)

Elles sont toutes au format \textsf{Ieeefloat\_64}, avec
\textsf{dimension = 8}, et \textsf{unit = second}
(le système unidata permettrait de rajouter  since 00h00 ).

Pour éviter de lire tout le fichier lorsque l'on recherche juste les heures de
début et de fin, on rajoute des champs \textsf{heure\_deb} et \textsf{heure\_fin},
au format \textsf{hh\string:mm\string:ss.ddd}

\section{Nommage}

  L'extension de nom utilisée pour ces fichiers est : \texttt{.vbn}.

  Le corps du nom de fichier (``\texttt{basename}'') correspond au
nom de la variable, tel qu'il est défini dans l'en-tête.

  Le nom d'un vecteur est à la fois un nom de fichier, et un nom de
variable, qui est utilisé pour faire référence à ces données (sélection
pour un graphique, formule de calcul, ...). La liste des caractères
possibles dans ce nom
cumule donc les contraintes des systèmes de fichiers et des interpréteurs.

\medbreak
  Les caractères suivants sont explicitement autorisés :

\begin{itemize}
	\item les lettres : de A à Z, en majuscule et en minuscul.
	\item les chiffres :  de 0 à 9
    \item le caractère souligné : « \_ »
    \item le caractère pourcent : « \% »
    \item le caractère 'deux points' : « : »
    \item le caractère 'deux points' : « : »
    \item le caractère arobase       : « @ »
    \item le caractère point         : « . »
\end{itemize}

\medbreak
  Il peut être envisagé d'utiliser les caractères suivants, qui ne sont pas
explicitement interdits  :
%– Enfin, il est recommandé que les programmes sachent manipuler des noms
%contenant les caractères qui n'ont pas encore été mentionnés («!#,:@[]
%{} ») mais le support de ces caractères n'est pas une obligation.

\begin{itemize}
  \item le caractère 'point d’exclamation' : « ! »
  \item le caractère 'dièze'               : « \# »
  \item le caractère 'virgule'             : « , »
  \item les caractères accolades ouvrante et fermante : « \{ » et « \} »
\end{itemize}

\medbreak
  A titre documentaire, les raisons d'exclusion des autres caractères sont
données en annexe, page \pageref{sec:exclu}

  D'un point de vue sécurité, l'interdiction de certains caractères ne veut
pas dire que les scripts ne doivent pas se protéger contre les noms de fichiers
qui contiendraient ces caractères. Cela signifie seulement qu'il est légitime,
dans le cadre AIDA, d'ignorer ou de signaler ces fichiers.

\textbf{Remarque} : les règles de nommage des paramètres (règles de
construction des noms en fonction de taxonomies et référentiels) ne sont pas
purement techniques, et sont donc décrites dans le volume 2 : manuel
d'utilisation (\url{user.pdf}).



% Partie : description des meta-données

\part{Organisation des meta-données}
\label{part:meta}
   Le terme de \textbf{meta-données} désigne les 
<<données sur les données>>, c'est à dire toutes
les informations sur les conditions de la mesure
et du traitement de celle-ci.

   Dans le cadre d'AIDA, on considère (principe d'auto-documentation)
que chaque vecteur contient l'information sur les 
conditions de sa propre production (attributs \textit{name},
\textit{unit}, \textit{origine},
\textit{fct\_origine}, \ldots). 

   Il reste donc à décrire : 
\begin{itemize}
  \item les informations propres à la configuration d'AIDA sur une machine
      particulière (liste et caractéristiques des plates-formes connues,
      liste des séries de mesures disponibles)
  \item des informations propres à l'architecture d'AIDA (par exemple
   la façon de regrouper les variables en familles).
  \item les informations relatives à une série de mesure (par exemple, 
      la description d'un vol, son découpage en paliers, \ldots) 
  \item les informations relatives aux campagnes de mesures
\end{itemize}

Il est important que toutes ces informations, même
celles qui évoluent peu, soient contenues dans des fichiers de
type texte, avec un format simple, car cette méthode est la
seule qui garantisse que les différents programmes, écrits dans
différents langages (C, Python, ...) partagent la même
information, et travaillent de façon cohérente.

Ces fichiers sont regroupés dans les répertoires : \\
\-\hspace{1cm} \path{${AIDA_META}} et \path{${AIDA_HOME}/config}. 

A l'exception notable du fichier donnant les détails
de chaque "acquisition" (ou période de mesure), tous
ces fichiers sont organisés de la même façon,
dans le format INI. Ces fichiers sont divisés en paragraphes débutant par un
nom écrit entre
crochets (par ex. \texttt{[merlin]}), suivi par des
lignes du type \texttt{attribut~=~valeur}. Les 
déclarations multi-lignes sont interdites, et les 
commentaires sont introduits par le caractère "\#"(dièse).

Ces fichiers n'ayant pas tous le même rôle, ils n'évolueront pas
de la même manière. Par exemple, ceux qui correspondent 
à l'organisation d'AIDA seront identiques sur
différents ordinateurs, et ne changeront que très
rarement. Ceux qui décrivent la liste des séries de
mesures disponibles évolueront rapidement, et pourront être
différents d'un site à l'autre.

\section{Description des plates-formes}
\label{sec:platef}
   La liste des plates-formes de mesure connues d'AIDA se trouve dans
le fichier \texttt{plates-formes}, dans le répertoire
\path{${AIDA_HOME}/profils/${AIDA_PROFIL}/}. Ce
fichier sert principalement à établir la relation entre le nom standard d'une
plate forme et l'abréviation qui permet de construire les
identificateurs uniques de chaque série de mesure 
(cf section \ref{sec:AcqId}), mais il contient aussi une
description un peu plus détaillée, qui est utilisée pour
renseigner les fichiers de documentation produits par
certains programmes de diffusion de données.

\subsection{Format}
  C'est un fichier au format "INI", divisé en paragraphes. Chaque
  paragraphe correspond à une plate-forme, et se présente comme suit :
\begin{description}
  \item[[Nom de la plate-forme]] tel qu'il est utilisé partout dans AIDA.
  \item[nom\_long] description en une ligne de texte, en français.
  \item[abrev] chaîne de deux caractères.
  \item[aliases] liste d'autres noms qui désignent parfois la plate-forme,
    séparés par des virgules.
  \item[nom\_serie\_mesure] nom d'une période de mesure dans la liste :
    Journee, Vol, Lacher, Lacher\_en\_vol, Drone  (la périodicité "Acquisition"
    n'est plus supportée et est équivalente à Journee)
  \item[duree\_vecteur] indique quelle est la durée d'un vecteur dans AIDA pour
    cette plate-forme. Les valeurs utilisées sont : "journee", "acquisition".
    Cette information est utilisée par les outils d'importation de données
    dans AIDA vers le format vecteur.
  \item[position] pour les plate-formes fixes, indique la position selon le
    format "latitude/longitude/altitude". Pour indiquer explicitement qu'une
    plate-forme n'est pas fixe, indiquer "mobile".\\
    Les coordonnées latitude et longitude sont écrites en degrés décimaux,
    positifs respectivement vers le Nord et l'Est.
\end{description}

  De manière optionnelle, et selon le type de plate-forme, des indicatifs
et/ou immatriculations peuvent être indiqués. Les syntaxes de ces champs sont :
\begin{description}
  \item[ICAO\_indicator] pour les plate-formes situées sur un aéroport,
    indicatif OACI de celui-ci.
  \item[WMO\_station\_index\_number] pour les stations synoptiques,
    indicatif OMM.
  \item[call-sign] pour les avions, immatriculation.
\end{description}

  Enfin, d'autres informations optionnelles peuvent être mentionnées :
\begin{description}
  \item[responsable] nom du reponsable de la plate-forme
\end{description}

\section{Description des acquisitions}
Pour chaque acquisition, il existe un fichier contenant toutes
ses méta-données.

Il est utilisable par les logiciels de visualisation et traitement,
ou simplement lisible par l'utilisateur. Il facilite le travail
de dépouillement des données.

Le format de ce fichier est décrit dans ce chapitre.

\subsection{Nom du fichier}

Pour une acquisition donnée le fichier ``inf'' s'appelle :
\texttt{\$\{Acq\_Id\}.inf}

 Acq\_Id est l'identificateur d'acquisition AIDA. Par exemple,
\path{fa110015} pour la quinzième acquisition de l'année 2011
sur le site du Fauga.
    Il n'y a aucune correspondance entre le numéro et la date de mesure.


\subsection{Stockage}

Tous les fichiers sont rangés dans le répertoire : 

\texttt{\$AIDA\_META\//INF}


\subsection{Format}

C'est un fichier de type texte, constitué d'une série de paragraphes,
eux-mêmes constitués d'une série d'égalités "champ = valeur". 

Le début d'un paragraphe est défini par la présence de son
nom en majuscules et entre crochets suivi d'un retour à la ligne.

A l'intérieur d'un paragraphe, la fin de ligne est marquée
par un point virgule:~\texttt{;}.

A l'intérieur de la valeur les informations sont séparées par un caractère:
\texttt{/}. Ces informations ne doivent pas contenir ce caractère \texttt{/}.

Les lignes vides, les espaces, et les tabulations sont permis.

Les lignes de commentaires doivent commencer par un caractère
dièse: \texttt{\#};

Les champs ne doivent pas contenir d'accents.

Les valeurs peuvent contenir des accents, à condition d'utiliser
le codage iso8859-1 (ou latin1).

Le fichier doit se terminer par \texttt{{[}FIN{]}}.

La liste des paragraphes ainsi que leur caractère "obligatoire"
ou "optionnel" est présentée dans le tableau ci-dessous.

\begin{center}

\begin{tabular}{|l|l|}
\hline 
Nom du Paragraphe& Caractère\\
\hline 
\hline 
\textbf{FORMAT}& \textbf{obligatoire}\\
\hline 
\textbf{IDENTIFICATION}& \textbf{obligatoire}\\
\hline 
\textbf{PARAMETRES\_ACQUIS}&\textbf{obligatoire}\\
\hline 
\textbf{ACQUISITION}& \textbf{obligatoire}\\
\hline 
HORAIRES\_UTC & optionnel\\
\hline 
INS& optionnel\\
\hline 
SEQUENCES& optionnel\\
\hline 
PERSONNEL& optionnel\\
\hline 
POINTS\_CARACTERISTIQUES& optionnel\\
\hline 
METEO\_DEBUT& optionnel\\
\hline  
METEO\_FIN& optionnel\\
\hline 
ETAT\_CAPTEURS& optionnel\\
\hline 
\end{tabular}
\end{center}


\subsubsection{{[}FORMAT{]}}

\label{parag-FORMAT}
Ce paragraphe décrit les formats utilisés dans les valeurs.

\ttfamily
\saute
\noindent\begin{tab2col}
\hline 
\textnormal{Champ}& \textnormal{Valeur}\\
\hline 
\hline 
format\_coord &
C (D)DDdMM.X(X) (C=lettre (N, E, S ou W), (D)DD=degrés 2 ou 3 chiffres, MM=minutes,
X(X)=fraction décimale de minutes;\\
\hline 
format\_heure&
HHMMSS (HH=heures, MM=minutes, SS=secondes);\\
\hline 
format\_date&
JJ/MM/AAAA (JJ=jour, MM=mois, AAAA=année, le tout en chiffres);\\
\hline 
format\_duree&
HH.D (HH=heure, D=dixièmes d'heure);\\
\hline 
\end{tab2col}
\saute

\normalfont

\subsubsection{{[}IDENTIFICATION{]}}
\label{parag-IDENTIFICATION}


\saute
\noindent\begin{tab4col}
\hline 
Champ& Caractère& Format valeur& Description \\
\hline 
\hline 
\texttt{plate-forme}& \textbf{obligatoire} & \textit{texte;}&
nom de la plate-forme en toutes lettres \\
\hline 
\texttt{acq\_id}& \textbf{obligatoire} & \textit{texte;} & identificateur d'acquisition AIDA\\
\hline 
\texttt{date}& \textbf{obligatoire} & \textit{JJ/MM/AAAA;} &
JJ= jour, MM= mois, AAAA= année (format décrit dans le paragraphe FORMAT)\\
\hline 
\texttt{campagne}& \textbf{obligatoire} & \textit{texte;}& nom de la campagne\\
\hline 
\texttt{theme}& optionnel& \textit{texte;}& thème scientifique de l'acquisition\\
\hline 
\texttt{lieu\_depart}& optionnel& \textit{texte;}&
lieu de départ de la plate-forme de mesure\\
\hline 
\texttt{lieu\_arrivee}& optionnel& \textit{texte;}&
lieu d'arrivée de la plate-forme de mesure\\
\hline 
\end{tab4col}
\saute

 Le nom de la plate-forme doit être connu d'AIDA par le fichier
\texttt{plate-forme}

Le champ \texttt{campagne} doit contenir une seule valeur.
Les acquisitions faites
hors de la période opérationnelle de la campagne ont comme nom de
campagne ESSAI, CALIB, ou TRANSIT.

Le champ \texttt{theme} doit expliciter plus librement le but de 
l'acquisition (ex : essais réfractomètre, vol type P1-flux, 
transit Brétigny-Dole, ....).

\subsubsection{{[}HORAIRES\_UTC{]}}

Ce paragraphe est utilisé dans le cas où la plate-forme de mesure est "mobile"
(avion, ballon, drone à ailes...)

\newcommand{\formathoraire}{\textit{hhmmss} ou\newline \textit{AAAA}-\textit{MM}-\textit{JJ}T\textit{hh}:\textit{mm}:\textit{ss}}
\saute
\noindent\begin{tab3col}
\hline 
Champ& Format valeur& Description\\
\hline 
\hline 
\texttt{depart\_parking}& \formathoraire & heure UTC de départ du parking\\
\hline 
\texttt{decollage}& \formathoraire & heure UTC de décollage\\
\hline 
\texttt{atterrissage}& \formathoraire & heure UTC d'atterrissage\\
\hline 
\texttt{retour\_parking}& \formathoraire & heure UTC de retour au parking\\
\hline 
\texttt{duree\_reglementaire}& \textit{hh.d}& durée réglementaire du vol\\
\hline 
\end{tab3col}
\saute


L'existence de deux formats impose aux programmes qui utilisent
ces données de savoir décoder les deux. La distinction se fait sur la
longueur du champ. La forme courte n'indique qu'une heure, et est donc ambigüe
dans les cas où l'acquisition s'étend sur plusieurs jours. Les différents cas
possibles sont :
\begin{itemize}
 \item utilisation d'un champ \textit{hh} >=24. Dans ce cas, l
 \item utilisation d'un champ \textit{hh} <24, mais avec une fin 
antérieure au début (par exemple \texttt{decollage = 223344} et 
\texttt{atterrissage = 001122})
\end{itemize}


Pour l'écriture de 
La forme longue du champ est recommandée, car elle n'est pas ambigüe,
et suit un format reconnu (ISO8601) par de nombreux langages et 
bibliothèques informatiques.

On reconnait que, lorsque les données sont saisies manuellement, 
la forme courte est plus concise et rapide. Il est cependant recommandé 
d'appliquer rapidement une normalisation après


\subsubsection{{[}ACQUISITION{]}}

\saute
\noindent\begin{tab4col}
\hline 
Champ& Caractère & Format valeur& Description \\
\hline 
\hline 
\texttt{systeme\_acquisition}& \textbf{obligatoire} & \textit{texte;}& nom du système d'acquisition\\
\hline 
\texttt{debut\_acquisition}& \textbf{obligatoire} & \textit{JJ/MM/AAAA/hhmmss;}& 
date et heure de début d'acquisition \\
\hline 
\texttt{fin\_acquisition}& \textbf{obligatoire} & \textit{JJ/MM/AAAA/hhmmss;}&
date et heure de fin d'acquisition \\
\hline 
\texttt{media}& \textbf{obligatoire} & \textit{texte;}& 
support de stockage du système d'acquisition\\
\hline 
\texttt{autre\_acquisition}& optionnel & \textit{texte/texte/...;}&
nom d'autres acquisitions sur la m\^eme plate-forme\\
\hline 
\texttt{commentaire}& optionnel & \textit{texte;}& commentaire\\
\hline 
\end{tab4col}
\saute

\noindent\texttt{media} peut prendre les valeurs floppy, dat, exabyte, hd,
don

\texttt{systeme\_acquisition} peut prendre les valeurs :
\begin{itemize}
\item RESUM, CL, PP, PPI pour l'ARAT
\item MC, MC95  pour le Merlin IV
\item EDISOL, ... pour les stations sols 
\item 4M-BATEAU, .... pour les acquisitions 4M  sur bateau
\item STAR, ... pour les radio-sondages
\end{itemize}

\subsubsection{{[}INS{]}}

Ce paragraphe donne les principales informations relatives à
la centrale à inertie de la plate-forme.

\saute
\noindent\begin{tab3col}
\hline 
Champ& Format valeur& Description\\
\hline 
\hline 
\texttt{ramp\_depart}& \textit{Nxxdxx.x/Wxxxdxx.x/xxxxft;}&
coordonnées du point de départ\\
\hline 
\texttt{type\_aln}&
\textit{xxx}~\texttt{\textit{s}}\textit{;} ou \textit{xxx}~\texttt{\textit{mn}}\textit{;}&
durée de l'alignement en secondes ou minutes\\
\hline 
\texttt{ramp\_arrivee}&
\textit{Nxxdxx.x/Wxxxdxx.x/xxxxft;}& coordonnées du point d'arrivée\\
\hline 
\texttt{pp\_arrivee}& \textit{Nxxdxx.x/Wxxxdxx.x/xxxxft;}&
coordonnées réelles à l'arrivée\\
\hline 
\end{tab3col}
\saute

\subsubsection{{[}POINTS\_CARACTERISTIQUES{]}}

Ce paragraphe décrit les différents points 
géographiques intéressant l'acquisition.

\saute
\noindent\begin{tab3col}
\hline 
Champ& Format valeur& Description\\
\hline 
\texttt{wpt}& \textit{nom/Nxxdxx.x/Wxxxdxx.x;}&
abréviation et coordonnées du point \\
\hline 
\texttt{radar}& \textit{nom/Nxxdxx.x/Wxxxdxx.x;}&
abréviation et coordonnées du point \\
\hline 
\texttt{balise}& \textit{nom/Nxxdxx.x/Wxxxdxx.x;}&
abréviation et coordonnées du point \\
\hline 
\texttt{bouee}& \textit{nom/Nxxdxx.x/Wxxxdxx.x;}&
abréviation et coordonnées du point \\
\hline 
\texttt{station}& \textit{nom/Nxxdxx.x/Wxxxdxx.x;}&
abréviation et coordonnées du point \\
\hline 
\texttt{ville}& \textit{nom/Nxxdxx.x/Wxxxdxx.x;}&
abréviation et coordonnées du point \\
\hline 
\end{tab3col}
\saute

\subsubsection{{[}PERSONNEL{]}}

\saute
\noindent\begin{tab3col}
\hline 
Champ&
Format valeur&
Description\\
\hline 
\hline 
\texttt{pilote}& \textit{Initiale du prénom. Nom;}&
Nom du pilote de la plate-forme (avion)\\
\hline 
\texttt{copilote}& idem&
Nom du copilote de la plate-forme\\
\hline 
\texttt{operateur}& idem&
Nom de l'opérateur du système d'acquisition\\
\hline 
\texttt{scientifique}& idem&
Nom du scientifique de la mission\\
\hline 
\texttt{directeur\_de\_mission}& idem&
Nom du directeur de la mission\\
\hline 
\texttt{invite}& idem&
Nom de l'invité \\
\hline 
\texttt{commandant\_de\_bord}& idem&
Nom du commandant de bord (bateau)\\
\hline 
\end{tab3col}
\saute

Les champs peuvent être répétés si plusieurs personnes
correspondent à la même désignation (ex: plusieurs 
scientifiques).


\subsubsection{{[}METEO\_DEBUT{]}}

Ce paragraphe décrit les conditions météorologiques 
en début d'acquisition.

\saute
\noindent\begin{tab3col}
\hline 
Champ& Format valeur& Description\\
\hline 
\hline 
\texttt{QNH}& \textit{xxx hPa;} ou \textit{xxx inHg;} 
ou  \textit{xxx mmHg;}& Valeur du QNH\\
\hline 
\texttt{QFE}& \textit{xxx hPa;} ou \textit{xxx inHg;} 
ou  \textit{xxx mmHg;} & Valeur du QFE\\
\hline 
\texttt{vent}& \textit{ddd/ff kt;} ou  \textit{ddd/ff m.s-1;}
& Direction en degrés et force du vent réel\\
\hline 
\texttt{temperature}& \textit{xxx C;} ou \textit{xxx F;} & Température en degrés C ou F\\
\hline 
\texttt{point\_de\_rosee}& \textit{xxx C;} ou \textit{xxx F;} & Température en degrés C ou F\\
\hline 
\texttt{visibilite}& \textit{xxxx m;} ou \textit{xx km;}& Visibilité\\
\hline 
\texttt{humidite}& \textit{xx \%;} & Humidité relative\\
\hline
\texttt{nuage}& \textit{texte;} & Description des nuages\\
\hline
\texttt{temps\_present}& \textit{texte;} & Description du temps présent\\
\hline
\texttt{origine}& \textit{texte;} & Origine de l'information météorologique\\
\hline
\texttt{commentaire}& \textit{texte;} & Commentaire de l'opérateur\\
\hline
\end{tab3col}
\saute

Le champ commentaire sert à recueillir les remarques de
l'opérateur sur l'information météorologique. Celles-ci
peuvent par exemple, décrire l'évolution rapide d'une situation,
ou apporter des informations non fournies dans les autres champs
du paragraphe.

\subsubsection{{[}METEO\_FIN{]}}

Ce paragraphe décrit les conditions météorologiques en fin
d'acquisition.

Le format est identique à \hbox{[METEO\_DEBUT]}.


\subsubsection{{[}SEQUENCES{]}}

Ce paragraphe décrit les différentes séquences spatiales et temporelles
au cours de l'acquisition.

\saute
\noindent\begin{tab2col}
\hline 
Champ& Format valeur\\
\hline 
\texttt{palier}&
\textit{nom/hhmmss\( _{deb} \)/hmmss\( _{fin} \)/Wpt\( _{deb} \)/Wpt\( _{fin} \)/Cap/Alt./comment;}\\
\hline 
\texttt{sondage}&
\textit{nom/hhmmss\( _{deb} \)/hmmss\( _{fin} \)/Wpt\( _{deb} \)/Wpt\( _{fin} \)/Alt\( _{deb} \)/Alt\( _{fin} \)/comment;}\\
\hline 
\texttt{cercle}&
\textit{nom/hhmmss\( _{deb} \)/hmmss\( _{fin} \)/Wpt/Alt/comment;}\\
\hline 
\texttt{zoom}&
\textit{nom/hhmmss\( _{deb} \)/hmmss\( _{fin} \)/Wpt\( _{deb} \)/Wpt\( _{fin} \)/Cap/Alt\( _{deb} \)/Alt\( _{fin} \)/comment;}\\
\hline 
\texttt{virage}&
\textit{nom/hhmmss\( _{deb} \)/hmmss\( _{fin} \)/Cap\( _{av} \)/Cap\( _{ap} \)/Angle parcouru/Alt/comment;}\\
\hline 
\texttt{station}&
\textit{nom/hhmmss\( _{deb} \)/hmmss\( _{fin} \)/Wpt/Cap/Alt/comment;}\\
\hline 
\texttt{sequence}&
\textit{nom/hhmmss\( _{deb} \)/hmmss\( _{fin} \)/comment;}\\
\hline 
\end{tab2col}
\saute

 \textit{nom} : abréviation de la séquence

 \textit{hhmmss\( _{deb} \)} : \textit{}heure de début de la
séquence sous la forme hhmmss

 \textit{hmmss\( _{fin} \)} : \textit{}heure de début de la
séquence sous la forme hhmmss

 \textit{Wpt} : abréviation du point géographique où se déroule
la séquence

 \textit{Wpt\( _{deb} \)} : abréviation du point de début de la
séquence

 \textit{Wpt\( _{fin} \)} : abréviation du point de fin de la
séquence

 \textit{Cap} : cap de la plate-forme durant la séquence en
degrés

 \textit{Cap$_{av}$} : cap de la plate-forme avant la séquence en
degrés 

 \textit{Cap$_{ap}$} : cap de la plate-forme après la séquence en
degrés

 \textit{Alt.} : altitude de la plate-forme durant la séquence
suivie de l'unité ft ou m, ou notée en niveau de vol FLxxx

 \textit{Alt\( _{deb} \)} : altitude de la plate-forme au début
de la séquence suivie de l'unité ft ou m, ou notée en niveau de vol FLxxx

 \textit{Alt\( _{fin} \)} : altitude de la plate-forme au début
de la séquence suivie de l'unité ft ou m,
ou notée en niveau de vol FLxxx

 \textit{Angle parcouru} : angle parcouru par la plate-forme au cours de la
séquence. Il est noté positif dans le sens des aiguilles d'une montre.

 \textit{comment} : commentaire texte


\subsubsection{{[}ETAT\_CAPTEURS{]}}

\saute
\noindent\begin{tab3col}
\hline 
Champ& Format valeur& Description\\
\hline 
\hline 
\texttt{avant\_acq}& \textit{texte;} &
commentaire sur l'état des capteurs avant l'acquisition\\
\hline 
\texttt{pendant\_acq}& \textit{texte;} &
commentaire sur l'état des capteurs pendant l'acquisition\\
\hline 
\texttt{apres\_acq}& \textit{texte;} &
commentaire sur l'état des capteurs après l'acquisition\\
\hline 
\end{tab3col}
\saute


\subsubsection{{[}PARAMETRES\_ACQUIS{]}}

\saute
\noindent\begin{tab3col}
\hline 
Champ& Format valeur& Description valeur\\
\hline 
\hline 
\texttt{abréviation}& \textit{texte;}& nom détaillé du paramètre\\
\hline 
\end{tab3col}
\saute

Ce paragraphe contient une ligne "\texttt{champ = valeur}" par paramètre
acquis. 
Le nom du champ est l'abréviation du paramètre.\\ 
Ce paragraphe peut être très long, car pour certaines plates-formes
le nombre de paramètres est important. Par conséquent pour améliorer la lisibilité
du fichier, ce paragraphe doit être \textbf{placé en fin de fichier}.





\section{Description des campagnes de mesures}
\label{sec:campagnes}
Une campagne de mesure est un ensemble de séries de mesures provenant de une ou plusieurs
plate-formes.

\subsection{Emplacement}
\texttt{\$AIDA\_META\//campagnes}
 
\subsection{Nommage}

Le fichier est nommé du nom de la campagne officiel transformé 
en majuscules avec l'extension \texttt{.camp}.\\
Dans le cas de campagnes ayant le même nom mais à quelques années d'intervalles on rajoutera 
l'année derrière le nom.

\subsection{Format}

C'est un fichier au format "INI", divisé en paragraphes.

Il contient un paragraphe IDENTIFICATION contenant les informations générales
de la campagne.
Il contient ensuite un paragraphe par plate-forme ayant participé à la campagne, nommé par le nom de la 
plate-forme défini dans le fichier plate-formes.

\subsubsection{[IDENTIFICATION]}

\saute
\noindent\begin{tabular}{|p{4.3cm}|p{2.2cm}|p{3.9cm}|p{5.5cm}|}
\hline 
Champ& Caractère& Format valeur& Description \\
\hline 
\hline 
\texttt{nom}& \textbf{obligatoire} & \textit{texte}& Nom de la campagne, identique au nom du fichier\\
\hline 
\texttt{date\_debut}& \textbf{recommandé} & \textit{JJ/MM/AAAA} & Date de début de la campagne\\
\hline 
\texttt{date\_fin }& \textbf{recommandé} & \textit{JJ/MM/AAAA} ou \textit{running} si la  campagne n'est pas terminée & Date de fin de la campagne \\
\hline 
\texttt{projet }& \textbf{recommandé} & \textit{texte} & Acronyme du projet. Le projet est plus large que la campagne. Il peut contenir plusieurs campagnes.\\
\hline 
\texttt{description }& \textbf{obligatoire} & \textit{texte} & Description succincte en anglais de la campagne et des objectifs scientifiques. \\
\hline
\texttt{region\_geographique }& \textbf{recommandé} & \textit{texte} & Texte en anglais. Exemple : South-West of France \\
\hline
\texttt{un\_m49\_code}& \textbf{recommandé} & \textit{Code un\_m49 de la région} & Codes disponibles sous : \href{https://en.wikipedia.org/wiki/UN_M49}{https://en.wikipedia.org/wiki/
UN\_M49}.\\
\hline 
\texttt{geospatial\_lat\_min }& \textbf{recommandé} & Degrés décimaux & Latitude minimale de la zone géographique explorée au cours de la campagne\\
\hline 
\texttt{geospatial\_lat\_max }& \textbf{recommandé} & Degrés décimaux & Latitude maximale de la zone géographique explorée \\
\hline 
\texttt{geospatial\_lon\_min }& \textbf{recommandé} & Degrés décimaux & Longitude minimale de la zone géographique explorée \\
\hline 
\texttt{geospatial\_lon\_max }& \textbf{recommandé} & Degrés décimaux & Longitude maximale de la zone géographique explorée \\
\hline 
\texttt{description }& \textbf{obligatoire} & \textit{texte} & Description succincte en anglais de la campagne et des objectifs scientifiques. \\
\hline 
\texttt{principal\_investigator }& \textbf{obligatoire} & \textit{prenom1:nom1:email1:
institution1; prenom2:nom2:
email2:institution2; ...} & PI(s) de la campagne globale au même format que celui utilisé dans nos fichiers NetCDF.\\
\hline 


\end{tabular}
\saute

\subsubsection{[<nom plate-forme>]}


\saute
\noindent\begin{tabular}{|p{4.3cm}|p{2.1cm}|p{4.7cm}|p{5cm}|}
\hline 
Champ& Caractère& Format valeur& Description \\
\hline 

\texttt{date\_debut}& \textbf{recommandé} & \textit{JJ/MM/AAAA} & Date de début de des mesures de la plate-forme dans la campagne\\
\hline 
\texttt{date\_fin }& \textbf{recommandé} & \textit{JJ/MM/AAAA} ou \textit{running} si la  campagne n'est pas terminée & Date de fin de des mesures de la plate-forme dans la campagne. \\
\hline 
\texttt{principal\_investigator }& \textbf{obligatoire} & \textit{prenom1:nom1:email1:
institution1; prenom2:nom2:
email2:institution2; ...} & PI(s) de la plateforme au même format que celui utilisé dans nos fichiers NetCDF. \\
\hline 

\texttt{liste\_acqid }& \textbf{obligatoire} & Liste d'Acq\_Id séparés par des virgules.\newline
Description d'une plage d'Acq\_Id sous la forme :\newline
\textit{<Acq\_Id\_deb>$-$<Acq\_Id\_fin>}
\textit{<Acq\_Id\_fin>} peut être égal à * si la campagne est en cours & Liste des identificateurs d'acquisition (Acq\_Id) correspondant à la campagne.\\
\hline 

\end{tabular}
\saute

\section{Description des familles de paramètres }
\label{sec:familles}

Il peut être utile de regrouper les paramètres en familles, ce qui
permet des tris thématiques, des aides à la saisie, etc.
Le fichier \path{familles}
décrit la liste des familles de paramètres connues d'AIDA.

\subsection{Emplacement}
 Le fichier \path{familles} est situé dans le répertoire
 \path{${AIDA_HOME}/config/}

\subsection{Format}
   C'est un fichier au format "INI". Il est encodé en
utf-8\footnote{contrairement à la plupart des fichiers d'AIDA,
qui sont codés en ISO8859-15}.
   Il est structuré sous la forme d'un paragraphe par famille. 

Chaque paragraphe se présente comme suit :
\begin{description}
  \item[[Nom du paragraphe]] nom de la famille, en codage ASCII sans accents (par exemple, \texttt{temperature})
  \item[acronyme = ]un groupe de trois lettres
  \item[nom:fr = ] le nom en français correct, avec accents (par exemple, \texttt{température})
  \item[nom:en = ] le nom en anglais
\end{description}
  S'il apparaît un besoin de traduction dans une autre langue,
il est possible de rajouter autant de lignes \textbf{\texttt{nom:<code\_langue>}} que nécessaire,
en utilisant pour code\_langue les valeurs de la table ISO639-1 (par
exemple \texttt{nom:ca} pour le catalan)


\section{Description des aéroports}
\label{sec:aeroports}
  Certains traitements ont besoin d'informations sur des aéroports : par
exemple, l'altitude de la piste, pour calculer une altitude-pression par
intégration à partir du décollage. Ces informations sont rassemblées dans
le fichier \path{airports.fdd}

\subsection{Emplacement}
  Le fichier \path{airports.fdd} est situé dans le répertoire
  \path{${AIDA_HOME}/profils/${AIDA_PROFIL}/}

\subsection{Format}
   C'est un fichier au format "INI", contenant un paragraphe par aéroport.
Le nom du paragraphe (entre crochets) est le nom de l'aéroport,
en français, sans accents, espace ni caractère de soulignement (\texttt{\_}).
Pour les noms composés, des tirets (\texttt{-}) seront
utilisés. Exemples : \texttt{La-Corogne}, \texttt{Baden-baden},
\texttt{Nimes-Courbessac}.

  Ce nom doit correspondre à celui utilisé pour les
champs \texttt{lieu\_depart} et \texttt{lieu\_arrivee}
du paragraphe \texttt{IDENTIFICATION} du fichier .inf.

  Les attributs possibles sont :\newline
\noindent\begin{tab3col}
\hline 
Champ& Format valeur& Description valeur\\
\hline 
\hline 
\texttt{ICAO\_ID}& \textit{4 caractères majuscule}& l'indicatif OACI de l'aéroport\\
\hline 
\texttt{Z0}& \textit{valeur numérique en mètres}& L'altitude de la piste (de l'endroit où se font la plupart des décollages)\\
\hline 
\end{tab3col}
\saute

\section{Description de l'évolution des paramètres de traitement.}
\label{sec:Evolution}
  Certains traitements utilisent des paramètres qui évoluent dans le temps. Ces
paramètres sont représentés par la syntaxe \$Chrono.PARAM dans la ligne de
commande adc. La valeur de PARAM est lue dans un fichier
\texttt{nom\_plateforme.chrono.}

\subsection{Emplacement}
  Les fichiers \path{nom_plateforme.chrono} sont situés dans le répertoire
  \path{${AIDA_HOME}/profils/${AIDA_PROFIL}/chrono}
 
\subsection{Format}
 Les lignes commençant par \# sont des commentaires.
 Les lignes blanches sont autorisées.
 \\
 
 Le fichier se compose d'une section par paramètre. La section est identifiée par le nom du paramètre entre crochets : [PARAM].
 Chaque section contient des lignes qui décrivent une plage de dates et la valeur de PARAM à l'intérieur de cette plage. Les bornes de la plage sont comprises. Le séparateur est constitué d'un ou plusieurs espaces ou tabulations.
 \\
 
 \textbf{Format des dates} : YYYY-MM-DDTHH:MI:SSZ
 
 \textbf{Valeurs possibles de PARAM} :
 
 -nombre (ex: -1.5)
        
 -chaîne (ex: "celsius")
      
 -nom de vecteur (ex: hum\_sol3cm\_ec5\_c3\_\%900)
 \\
 
 
 \textbf{Exemple de fichier chrono} : 
 
 \begin{lstlisting}
 # Chrono toulouse-flux
 
 # délai du LICOR
 [DELAI]
 2012-06-15T00:00:00Z    2014-10-14T23:59:59Z    -0.200
 2014-10-15T00:00:00Z    2015-05-28T23:59:59Z    -0.250
 2015-05-29T00:00:00Z    2020-12-31T23:59:59Z    -0.200
 
 # unité températures pour tests
 [UNITE_TPR]
 2012-06-15T00:00:00Z    2014-07-09T23:59:59Z    "Kelvin"
 2014-07-10T00:00:00Z    2020-12-31T23:59:59Z    "Celsius"
 
 # vecteur pour tests
 [VECTEUR]
 2014-06-15T00:00:00Z    2014-07-09T23:59:59Z    hum_sol3cm_ec5_c3_%900

 2014-07-10T00:00:00Z    2020-12-31T23:59:59Z    hum_sol3cm_ec5_c3_bis_%900
 \end{lstlisting}

 \saute 
 
 


% Annexes

% Annexes
\part{Annexes}
\section{exclusion des autres caractères}
\label{sec:exclu}
\noindent

\begin{tabular}{p{0.33\textwidth}p{0.66\textwidth}}
  Caractère(s) & Raison\\ 
  \hline
  Caractère hors du jeu ASCII 7~bits 
  et l'espace.
  & Même si les
  systèmes de fichier des systèmes d’exploitation actuels savent gérer accents
  , espaces et caractères non-latins, il subsiste des problèmes de portabilité
  (notamment lors de l'utilisation de NFS ou SMB entre systèmes avec 
  des codages différents), et il n'y
  a pas suffisamment d'avantages à utiliser les accents pour accepter les
  inconvénients associés.\\
  \hline
  Les caractères de contrôle, de code entre 0 et
  31, et 127(« del ») & Ne sont pas faits pour cet usage. \\
  \hline
  \texttt{/} (barre oblique) et\newline
  \texttt{\textbackslash} (barre oblique inversée) &
   Ils servent de séparateurs de nom de répertoire dans les
  chemins d'accès aux fichiers sur différents systèmes d'exploitation et dans
  les URL/URI.\\
  \hline
  \texttt{"} (guillemet droit),\newline
  \texttt{*} (astérisque ou étoile),\newline
  \texttt{?} (point d'interrogation),\newline
  \texttt{>} et \texttt{<} (supérieur et inférieur),\newline
  \texttt{|} (barre verticale)
   & Ils sont impossibles sur certains systèmes de fichiers.\\
  \hline
  \texttt{'} (apostrophe),\newline
  \texttt{\$} (dollar),\newline
  \texttt{;} (point-virgule),\newline 
  \texttt{\&} (esperluette) et\newline
  \texttt{`} (guillemet ouvrant simple) & Les
  inconvénients qu'ils engendrent pour les utiliser correctement dans des
  environnements usuels (shell unix, par exemple) ne sont pas supérieurs à
  l'intérêt que pourrait avoir leur utilisation\\
  \hline
  \texttt{+} (plus),\newline
  \texttt{-} (moins, trait d'union ou tiret),\newline
  \texttt{\^{}} (accent circonflexe avec chasse),\newline
  \texttt{(} et \texttt{)} (parenthèses ouvrante et fermante),\newline
  \texttt{[} et \texttt{]} (crochets ouvrant et fermant) et\newline
  \texttt{=} (signe égal) & Ils servent aux opérations
  mathématiques usuelles lors de l'interprétation d'une formule (avec « * » et
  « / ». qui sont déjà impossibles pour d'autres raisons).\\
  \hline
\end{tabular}


\end{document}
